\subsection{EXPONENTIAL AND LOGARITHM IN FILTERED ALGEBRAS}

Let \(\mathbf{A}\) be a unital associative algebra which is Hausdorff and complete under a real filtration \((\mathbf{A}_{\alpha})\) . We write \(m = A_0^+ = \bigcup_{\alpha > 0} A_\alpha\) .

For \(x \in \mathfrak{m}\) , the family \((x^n / n!)_{n \in \mathbf{N}}\) is summable. We write

\[
e ^ {x} = \exp x = \sum_ {n \geqslant 0} x ^ {n} / n!.\tag{1}
\]

Then \(\exp (x)\in 1 + \mathfrak{m}\) and the mapping \(\exp \colon \mathfrak{m}\to 1 + \mathfrak{m}\) is called the exponential mapping of A.

For all \(y \in 1 + \mathfrak{m}\) , the family \(((-1)^{n-1}(y - 1)^n / n)_{n \geqslant 1}\) is summable. We write

\[
\log y = \sum_ {n \geqslant 1} (- 1) ^ {n - 1} (y - 1) ^ {n} / n.\tag{2}
\]

Then \(\log y \in m\) and the mapping \(\log: 1 + m \to m\) is called the logarithmic mapping of A.

PROPOSITION 1. The exponential mapping is a homeomorphism of \(\mathfrak{m}\) onto \(1 + \mathfrak{m}\) and the logarithmic mapping is the inverse homeomorphism.

For \(x \in \mathrm{A}_{\alpha}, \frac{x^n}{n!} \in \mathrm{A}_{n\alpha}\) . It follows that the series defining the exponential converges uniformly on each of the sets \(\mathrm{A}_{\alpha}\) for \(\alpha > 0\) ; as \(\mathrm{A}_{\alpha}\) is open in \(m\) and \(m = \bigcup_{\alpha > 0} \mathrm{A}_{\alpha}\) , the exponential mapping is continuous. It can be similarly shown that the logarithmic mapping is continuous.

Let \(e\) and \(l\) be the formal power series with no constant term

\[
e (\mathrm{X}) = \sum_ {n \geqslant 1} \frac {\mathrm{X} ^ {n}}{n !}, \quad l (\mathrm{X}) = \sum_ {n \geqslant 1} (- 1) ^ {n - 1} \mathrm{X} ^ {n} / n.
\]

We know (Algebra, Chapter IV, § 6, no. 9) that \(e(l(\mathbf{X})) = l(e(\mathbf{X})) = \mathbf{X}\) on \(\hat{\mathrm{A}}(\{\mathbf{X}\}) = \mathrm{K}[[\mathbf{X}]]\) . By substitution (§ 5, no. 1), we deduce that

\[
e (l (x)) = l (e (x)) = x
\]

for \(x\in \mathfrak{m}\) ; as

\[
\exp x = e (x) + 1, \quad \log (1 + x) = l (x)
\]

it follows immediately that

\[
\log \exp x = x, \quad \exp \log (1 + x) = 1 + x
\]

for \(x\) in \(\mathfrak{m}\) , whence the proposition.

Remarks. (1) If \(x \in \mathfrak{m}\) , \(y \in \mathfrak{m}\) and \(x\) and \(y\) commute, then

\[
\exp (x + y) = \exp (x) \exp (y),
\]

since the family \(\left(\frac{x^i}{i!}\cdot \frac{y^j}{j!}\right)_{i,j\in \mathbf{N}}\) is summable (cf.~Algebra, Chapter IV, § 6, no. 9, Proposition 11).

\begin{enumerate}
\def\labelenumi{(\arabic{enumi})}
\setcounter{enumi}{1}
\item
  As the series \(e\) and \(l\) are without constant term and \(A_{\alpha}\) is a closed ideal of \(A\) , \(\exp A_{\alpha} \subset 1 + A_{\alpha}\) and \(\log (1 + A_{\alpha}) \subset A_{\alpha}\) whence \(\exp A_{\alpha} = 1 + A_{\alpha}\) and \(\log (1 + A_{\alpha}) = A_{\alpha}\) for \(\alpha > 0\) .
\item
  Let B be a complete Hausdorff filtered unital associative algebra and \(n = \bigcup_{\alpha > 0} B_{\alpha}\) . Let f be a continuous unital homomorphism of A into B such that \(f(\mathfrak{m}) \subset \mathfrak{n}\) . Then \(f(\exp x) = \exp f(x)\) for \(x \in \mathfrak{m}\) and \(f(\log y) = \log f(y)\) for \(y \in 1 + \mathfrak{m}\) ; we show for example the first of these formulae:
\end{enumerate}

\[
f (\exp x) = \sum_ {n \geqslant 0} f (x ^ {n}) / n! = \sum_ {n \geqslant 0} f (x) ^ {n} / n! = \exp f (x).
\]

\begin{enumerate}
\def\labelenumi{(\arabic{enumi})}
\setcounter{enumi}{3}
\tightlist
\item
  Let E be a unital associative algebra. If \(a\) is a nilpotent element of E, the family \(\left(\frac{a^n}{n!}\right)_{n\in \mathbf{N}}\) has finite support and we write \(\exp a = \sum_{n\geqslant 0}a^n /n!\) . An element \(b\) is unipotent if \(b - 1\) is nilpotent; then we write
\end{enumerate}

\[
\log b = \sum_ {n \geqslant 1} (- 1) ^ {n - 1} (b - 1) ^ {n} / n.
\]

We deduce from the relations \(e(l(\mathbf{X})) = l(e(\mathbf{X})) = \mathbf{X}\) that the mapping \(a \mapsto \exp a\) is a bijection of the set of nilpotent elements of E onto the set of unipotent elements of E and that \(b \mapsto \log b\) is the inverse mapping.

